\documentclass[11pt]{article}

\usepackage[a4paper,margin=28mm]{geometry}
\usepackage{amsmath,amssymb,amsthm,mathtools}
\usepackage{booktabs}
\usepackage{enumitem}
\usepackage[hidelinks]{hyperref}
\usepackage{xcolor}

\hypersetup{
  pdftitle={Prime-support restrictions for a magic square of distinct squares},
  pdfauthor={Anonymous AI-generated research draft},
  pdfsubject={Unverified AI-generated partial results on the magic square of squares}
}

\newtheorem{theorem}{Theorem}[section]
\newtheorem{proposition}[theorem]{Proposition}
\newtheorem{lemma}[theorem]{Lemma}
\newtheorem{corollary}[theorem]{Corollary}
\theoremstyle{definition}
\newtheorem{definition}[theorem]{Definition}
\theoremstyle{remark}
\newtheorem{remark}[theorem]{Remark}

\newcommand{\Z}{\mathbb Z}
\newcommand{\Q}{\mathbb Q}
\newcommand{\F}{\mathbb F}
\newcommand{\Norm}{\operatorname N}

\title{Prime-support restrictions for a magic square of distinct squares\\
\large An AI-generated research report and verification index}
\author{Anonymous AI-generated research draft}
\date{20 August 2026}

\begin{document}

\maketitle

\begin{abstract}
This AI-generated research project studies the open question of whether a
$3\times3$ magic square can have nine pairwise distinct positive integer
squares.  We study the prime support of
the square root $e$ of the center entry $e^2$.  The main results recorded here
exclude every primitive center root supported on exactly two rational primes,
with arbitrary positive exponents, and exclude every exactly-three-block
center for which all local Gaussian indices are balanced or extreme.  The
latter includes every squarefree center root $pqr$.  For the first remaining
exponent pattern $p^2qr$, an exact physical signature classification gives 134
classes; proved valuation, quotient, factorization, and congruence filters
exclude 118 and leave 16 necessary classes.  No surviving class is asserted
to be realizable.  We also describe bounded exact searches, used only as
cross-checks.  The unrestricted problem is not solved.

This report is a compact statement and proof index.  A versioned companion
archive contains the line-by-line Gaussian calculations, elliptic-curve
descents, independent finite enumerators, scanners, and regression tests.
The work has not been peer reviewed and makes no priority claim.
\end{abstract}

\begin{center}
\fcolorbox{black}{gray!8}{\parbox{0.92\textwidth}{\small
\textbf{AI-generated mathematical result.}
AI systems generated most of the conjectures, proof strategies, derivations,
code, and exposition under human direction.  Uwe Schwarz contributed the
problem choice, prompts, computing resources, coordination, and review; he
does not claim authorship of the mathematical results.  Repeated AI review and
exact tests are not independent expert verification.  The work is public to
make errors discoverable and has not been peer reviewed.
}}
\end{center}

\tableofcontents

\section{Problem, conventions, and scope}

A normal $3\times3$ magic square has equal row, column, and diagonal sums.
The open problem asks whether all nine entries can be distinct positive
integer squares.  Since the center of every $3\times3$ magic square is one
third of the common line sum, a square-valued center may be written $e^2$.
Dividing the nine square roots by their common greatest common divisor reduces
the problem to a \emph{primitive} square without changing distinctness.

This report proves restrictions on the factorization of $e$.  It does not
prove existence or global nonexistence.  In particular, the following remain
open here:
\begin{itemize}[leftmargin=2em]
  \item centers with four or more rational prime factors;
  \item exactly-three-block centers with genuine intermediate local indices;
  \item the 16 weighted signatures left by the $p^2qr$ analysis;
  \item construction of a full square, or a global obstruction to all centers.
\end{itemize}

The literature still describes the unrestricted problem as open.  Bremner
gave elliptic-curve and intersection-of-quadrics formulations and strong
near-solutions \cite{bremner1999,bremner2001}.  More recent work studies
Gaussian and residue reformulations \cite{woll2018,cain2019,wolird2023}, while
Rome and Yamagishi settle the existence problem for orders at least four, not
order three \cite{rome2025}.  A bounded literature audit and discussion of
recent claims is included in the companion archive.

\section{Centered form and four offsets}

Every $3\times3$ magic square with center $e^2$ has the form
\begin{equation}
\begin{pmatrix}
e^2+x & e^2-x-y & e^2+y\\
e^2-x+y & e^2   & e^2+x-y\\
e^2-y & e^2+x+y & e^2-x
\end{pmatrix}.
\label{eq:centered}
\end{equation}
After interchanging the two centered coordinates and changing their signs,
put $a>b>0$.  The four positive offsets of opposite pairs are
\begin{equation}
\mathcal D=\{a,b,a+b,a-b\}.
\label{eq:offsets}
\end{equation}
The entries are pairwise distinct exactly when these four offsets are nonzero
and pairwise distinct.  For every $d\in\mathcal D$ there are positive integers
$u_d,v_d$ satisfying
\begin{equation}
u_d^2=e^2+d,\qquad v_d^2=e^2-d,
\end{equation}
and therefore
\begin{equation}
(u_dv_d)^2+(d)^2=e^4.
\label{eq:circle}
\end{equation}
Thus each offset is the absolute imaginary coordinate of a Gaussian integer
square of norm $e^4$.  The two corner offsets $a,b$ and two edge offsets
$a+b,a-b$ also give four three-term relations.  Deleting an edge gives a
coefficient-$111$ relation; deleting a corner gives a coefficient-$211$
relation.

\section{Gaussian blocks and local indices}

Let $p^k\Vert e$.  A primitive representation of \eqref{eq:circle} forces
$p\equiv1\pmod4$, so choose $p=\pi\bar\pi$ in $\Z[i]$.  At the $p$-support,
each offset has an index
\begin{equation}
j_{p,d}\in\{-k,-k+1,\ldots,k-1,k\}.
\end{equation}
The values $j=0$ and $|j|=k$ will be called \emph{balanced} and
\emph{extreme}.  Values $0<|j|<k$ are \emph{intermediate}.  Unique
factorization in $\Z[i]$ gives the following local spectrum.

\begin{lemma}[Local unit pattern]
For every center prime $p$, at least three of the four offsets are $p$-adic
units.  Equivalently, at most one column has a nonextreme local valuation
pattern.  The possible exception is constrained by whether it is a corner or
an edge and by the coefficient-$111$ or coefficient-$211$ relation obtained
from the other three columns.
\end{lemma}

The proof uses the pairwise differences among $a,b,a+b,a-b$.  If two offsets
had positive $p$-valuation, then the displayed linear relations force a third
and finally contradict primitivity.  Once three offsets are units, aligning
their one-sided Gaussian factors produces an exact rational quotient.  These
quotients are the main bridge between local valuations and global sizes.

\section{A block-balance theorem}

Write $M=e/p^k$.  The central quantitative result is the following.

\begin{theorem}[Prime-power block balance]
\label{thm:block}
For every $p^k\Vert e$ in a primitive magic square of distinct squares,
\begin{equation}
p^{2k}<3M^2.
\label{eq:sqrt3}
\end{equation}
If $p\equiv5\pmod8$, then the stronger inequality
\begin{equation}
p^k<M
\label{eq:mod8}
\end{equation}
holds.  If $p^k\ge M$, there is exactly one nonunit offset; it is an edge,
$p\equiv1\pmod8$, and for its intermediate index $j$,
\begin{equation}
k\ge3,\qquad
\left\lfloor\frac{k}{2}\right\rfloor+1\le |j|\le k-1,
\qquad p^{2(k-|j|)}<\sqrt2 M.
\label{eq:dominant}
\end{equation}
\end{theorem}

\paragraph{Proof architecture.}
Align the three $p$-unit Gaussian squares in the retained $111$ or $211$
relation and remove their common one-sided factor.  Unique factorization gives
an identity
\begin{equation}
H=\bar\Pi n,\qquad n\in\Z\setminus\{0\},
\end{equation}
or the same identity for $H/2$ in the $211$ case.  All residual summands have
the same modulus.  Strict triangle inequalities give $|n|p^{2k}<3M^2$ in the
$111$ case and the constant 2 in the $211$ case.  The valuation spectrum shows
that the factor represented by $n$ contains precisely the missing local
content.  If a block is at least as large as its complement, the corner case
has incompatible $p$-valuation, while an edge case survives only with
$(2/p)=1$.  This proves \eqref{eq:sqrt3}--\eqref{eq:dominant}.  The detailed
unit choices, corner contents, and Legendre-symbol refinement are given in
\texttt{research/prime-support/block-balance.md}.

An immediate consequence is that a center cannot be a prime power.  More
importantly, \eqref{eq:sqrt3} can be applied simultaneously to every block;
the resulting strict monomial inequalities are strong enough to eliminate
many complete support patterns.

\section{Centers supported on two primes}

\begin{theorem}[Two-block exclusion]
\label{thm:two}
There is no primitive magic square of distinct positive integer squares whose
center root is
\begin{equation}
e=p^kq^\ell
\end{equation}
for distinct rational primes $p,q$ and positive integers $k,\ell$.
\end{theorem}

The proof is arithmetic, not a bounded search.  Its main reductions are:
\begin{enumerate}[leftmargin=2em]
  \item neither block can have a balanced exceptional offset;
  \item the two exceptions must both be edges;
  \item after ordering $P=p^k>Q=q^\ell$, the larger block is dominant, so
        $k\ge3$ and $p\equiv1\pmod8$;
  \item the two edge relations share a forced sign parameter and a small odd
        quotient $\nu$;
  \item congruences and the coupled norm equations reduce to $|\nu|=3$ and a
        genus-one curve.
\end{enumerate}
In the final branch, rational parameters lead to
\begin{equation}
y^2=t^4-10t^2+1.
\label{eq:quartic-two}
\end{equation}
A birational transformation gives
\begin{equation}
E:\quad v^2=x(x-1)(x-3).
\label{eq:curve-two}
\end{equation}
A self-contained $2$-isogeny descent shows that $E(\Q)$ has rank zero; good
reduction at 5 and 7 bounds the torsion, and the visible points exhaust it.
Every torsion point maps back to $t=0$ or to a degenerate offset.  This closes
the last edge-edge branch and proves Theorem~\ref{thm:two}.  The complete sign
table, norm elimination, covering equations, local obstructions, and inverse
maps appear in
\texttt{research/prime-support/two-block-exponent-one.md}.

\begin{remark}
The arbitrary exponents in Theorem~\ref{thm:two} matter.  A squarefree
semiprime exclusion alone does not imply the theorem.
\end{remark}

\section{Exactly three blocks in the balanced/extreme regime}

Let
\begin{equation}
e=PQR=p^kq^\ell r^m
\end{equation}
have exactly three prime-power blocks.  Suppose every local index is balanced
or extreme.  After normalizing the three one-sided phases, a column is encoded
by a vector in $\{0,\pm1\}^3$ up to global sign.  Distinct offsets require the
four projective vectors to be nonzero and pairwise distinct.  Every deleted
triple must also be orientation-frustrated: a coherent $111$ or $211$ triple
would give a rational Gaussian quotient whose one-sided prime supports cannot
match.

The physical symmetry group is smaller than an arbitrary permutation of four
columns: it preserves the pair of corners and the pair of edges.  Quotienting
by row permutations, corner swap, edge swap, row switches, and column
conjugations leaves a finite list of physical signatures.

\begin{theorem}[Three balanced/extreme blocks]
\label{thm:three}
No primitive center root supported on exactly three rational primes is
admissible if, at every block, all four local indices are balanced or extreme.
\end{theorem}

\begin{corollary}[Squarefree three-prime exclusion]
There is no primitive magic square whose center root is $e=pqr$ for three
distinct rational primes.
\end{corollary}

\paragraph{Finite and arithmetic parts.}
In the squarefree model, exact projective distinctness and frustration leave
29 physical classes, including one all-extreme Hadamard class.  The
all-extreme class reduces to two elliptic curves isomorphic to
\begin{equation}
Y^2=X(X+1)(X+5),
\end{equation}
whose rational points are only the degenerate torsion points.  Thus 28 classes
remain.  They are excluded by a sequence of exact arguments: elementary
trigonometric inequalities, Fermat's right-triangle descent, simultaneous
square lemmas, Gaussian quotient comparisons, and three rank-zero elliptic
curves.  None of these 28 exclusions is a bounded search.

For arbitrary exponents, replace each prime by its full block size and each
one-sided Gaussian prime by the corresponding power.  Primitivity,
non-axis conditions, projective distinctness, quotient integrality, and every
descent survive unchanged.  This transfers all 28 exclusions and proves
Theorem~\ref{thm:three}.  The classifiers and proof are in
\texttt{three-block-squarefree.md},
\texttt{three\_block\_signatures.py}, and
\texttt{three\_block\_squarefree.py}.

\section{The first intermediate pattern: $p^2qr$}

Theorem~\ref{thm:three} forces any $p^2qr$ survivor to have one intermediate
$p$-index $\pm1$; its other $p$-indices are $\pm2$, while the $q$- and
$r$-rows are balanced or extreme.  Exact physical quotienting gives 134
weighted signature classes.

\begin{proposition}[Verified $p^2qr$ residue]
\label{prop:p2qr}
Of the 134 necessary physical signatures for $e=p^2qr$, the proved filters in
the companion archive exclude 118.  Exactly 16 classes remain relative to
those filters.  No remaining class is asserted realizable.
\end{proposition}

The reduction is intentionally staged so that every count is independently
testable.
\begin{center}
\begin{tabular}{@{}lr@{}}
\toprule
Stage & Total classes excluded \\
\midrule
Exceptional-row monomial inequalities & 58 \\
Aligned rational quotients & 85 \\
Composite-block bounds & 95 \\
General common-factor heights & 104 \\
All admissible divisor comparisons & 110 \\
Factored axis bounds & 113 \\
Congruence-sharpened skew bounds & 116 \\
Discrete prime floors & 117 \\
Complementary-axis descent & 118 \\
\bottomrule
\end{tabular}
\end{center}

The filters combine strict monomial inequalities with exact local data.  For
example, if an aligned retained triple has a common Gaussian factor, rational
valuation symmetry turns its residual quotient into an integer.  Factoring
the remaining axis expression gives divisibility by one cyclotomic factor,
while skew expressions produce a norm quotient with a congruence gap.  The
sharp uniform constants used in the final classifier are
\begin{equation}
S_2=241,\qquad S_6=433,\qquad S_4=1201,\qquad S_8=4201.
\end{equation}
All cone comparisons use exact integer arithmetic, including constants; no
floating-point feasibility decision is made.  The 134-class enumerator and
the arithmetic-filter implementation are independent programs.  Their tests
pin every intermediate count in the table.

The 16 survivors are a concrete frontier rather than evidence of existence.
They retain enough sign and intermediate-index freedom that the current local
quotients do not close.  A productive next step would be to couple two or more
of their exact Gaussian identities rather than add another independent height
bound.

\section{Bounded exact computation}

The proof results above do not depend on search limits.  Separate scanners
were used to find counterexamples to intermediate lemmas and to cross-check
enumerators.  Their recorded negative results are:
\begin{center}
\begin{tabular}{@{}p{0.43\textwidth}p{0.36\textwidth}r@{}}
\toprule
Search & Exhausted range & Hits \\
\midrule
Center-root $111$ and $211$ relations & $e\le 5\cdot10^8$ & 0 \\
Rational-parameter relation scan & parameter height $H\le5.5\cdot10^6$ & 0 \\
Dominant-block identity scan & complementary block $M\le10^6$ & 0 \\
Two-block factor scan & $\max(P,Q)\le10^{12}$ & 0 \\
\bottomrule
\end{tabular}
\end{center}
The largest rational-parameter run used four threads and enumerated 875,336
normalized values.  It tested exactly $383{,}106{,}118{,}780$ unordered pairs
and completed in 8 hours, 42 minutes on the host used for the scan.  The
center-root and factor scans use different heights, so their zero counts are
not interchangeable.  This larger rational-parameter run is useful as
evidence but cannot settle the Diophantine question.

An independent B6 project reports no full nine-entry candidate through its
root-box bound $R\le10^6$; see the project's
\href{https://github.com/mystimath/magic-square-of-squares-3x3/blob/main/docs/40-b4-b6-publication-revalidation-2026-07-21.md}{revalidation report}.
That external computation was source-checked here but not reproduced.  Its
root-box height is not interchangeable with any of the local bounds above.

Every scanner has synthetic positive tests and a small-range independent
reference enumerator.  The repository's regression suite currently contains
46 tests covering the finite classifications, arithmetic stages, Gaussian
arithmetic, native scanner builds, and reference comparisons.

\section{Reproducibility and proof map}

From the repository root, run
\begin{verbatim}
make test
make paper
\end{verbatim}
The first command executes the Python regression suite; several tests compile
the C++20 scanners in temporary directories.  The second command uses Tectonic
to rebuild this report.

The principal proof files are:
\begin{center}
\begin{tabular}{@{}p{0.35\textwidth}p{0.55\textwidth}@{}}
\toprule
File & Content \\
\midrule
\texttt{block-balance.md} & local valuation spectrum, block inequalities,
dominant exceptions, Legendre filters \\
\texttt{two-block-exponent-one.md} & arbitrary-exponent two-block exclusion
and final $2$-isogeny descent \\
\texttt{hourglass-sunit.md} & primitive denominators and frustration of
$111$/$211$ relations \\
\texttt{composite-blocks.md} & compatible orientation grouping and general
common-factor bounds \\
\texttt{three-block-squarefree.md} & 28 physical classes, arithmetic
exclusions, and arbitrary-exponent transfer \\
\texttt{three-block-p2qr.md} & 134-class weighted residue and the complete
118-stage filter proof \\
\texttt{literature-audit.md} & source status, overlap, and novelty cautions \\
\bottomrule
\end{tabular}
\end{center}

This separation is deliberate.  The report makes the scope and dependency
graph readable; the companion notes preserve enough algebra to audit signs,
units, valuations, exceptional cases, and elliptic-curve maps line by line.

\section{Assessment and open review questions}

The results appear materially stronger than a routine first computation.
The two-block theorem allows arbitrary exponents; the three-block theorem
transfers an exact physical classification and several descents beyond the
squarefree case; and the $p^2qr$ work exposes a finite, reproducible frontier.
At the same time, no systematic MathSciNet or zbMATH priority search has been
performed, and the methods overlap with established Gaussian-integer and
elliptic-curve approaches.  The appropriate public claim is therefore
``independently checkable partial progress,'' not ``new solution'' and not a
priority assertion.

The most useful independent reviews would address:
\begin{enumerate}[leftmargin=2em]
  \item the unit and sign alignment in the coupled two-block edge relations;
  \item the transfer from squarefree to arbitrary balanced/extreme blocks;
  \item completeness of the physical symmetry quotient in the three-block
        classifiers;
  \item the local-to-global step in each common-factor quotient;
  \item the elliptic-curve descents and exceptional inverse images;
  \item whether any of the 16 $p^2qr$ survivors can be eliminated by a coupled
        identity or realized locally at all primes.
\end{enumerate}

\section*{Acknowledgment}

The public problem page at UnsolvedMath and the sources below provide context
for independent discussion.  Uwe Schwarz contributed human direction,
computing resources, curation, and review without claiming mathematical
authorship.  The project actively invites independent verification and
correction.

\begin{thebibliography}{9}

\bibitem{bremner1999}
A. Bremner,
\emph{On squares of squares},
Acta Arithmetica 88 (1999), 289--297.
\url{https://eudml.org/doc/207247}

\bibitem{bremner2001}
A. Bremner,
\emph{On squares of squares II},
Acta Arithmetica 99 (2001), 289--308.
\url{https://doi.org/10.4064/aa99-3-6}

\bibitem{woll2018}
C. Woll,
\emph{A Partial Residue Categorization of the Magic Square of Squares},
arXiv:1809.03067 (2018).
\url{https://arxiv.org/abs/1809.03067}

\bibitem{cain2019}
O. Cain,
\emph{Gaussian Integers, Rings, Finite Fields, and the Magic Square of
Squares}, arXiv:1908.03236 (2019).
\url{https://arxiv.org/abs/1908.03236}

\bibitem{wolird2023}
C. Wolird,
\emph{A New Transformation of the Magic Square of Squares},
arXiv:2310.12164 (2023).
\url{https://arxiv.org/abs/2310.12164}

\bibitem{rome2025}
N. Rome and S. Yamagishi,
\emph{On the existence of magic squares of powers},
Mathematika 71 (2025), e70016.
\url{https://arxiv.org/abs/2406.09364}

\end{thebibliography}

\end{document}
